\scanpage{103}
elle \unsure{vaut} sur $1,U_1,U_2$ par
\[
1\longmapsto2\deleted{\unsure{e}},\qquad U_1\mapsto-b_1,\quad U_2\mapsto-b_2
\]
\unsure{chute} en $\Gamma\underline{\mathcal O}\subset\Gamma\mathcal A$, il reste à trouver l'image de $U_1\otimes U_2$, c'est donc un élément de $\mathcal A$ dont l'image dans $(\mathcal A_1\otimes\mathcal A_2)^{\unsure{\natural}}$ est
\[
U_1\otimes U_2+\underbrace{\bar U_1\otimes\bar U_2}_{(-b_1\unsure{-|+}U_1)\otimes(-b_2-U_2)}=\underbrace{2U_1U_2+b_1U_2+b_2U_1}_{\text{image de }U}+b_1b_2
\]
\deleted{\unsure{donc}} donc il y a lieu de prendre
\[
U_1*U_2=U+b_1b_2
\]
Mais il faut maintenant définir une application
\[
N_{\mathcal A_1\otimes\mathcal A_2/\mathcal A}:\mathcal A_1\otimes\mathcal A_2\unsure{\mathrel{-\!\!-\!\!-}|\longrightarrow}\mathcal A
\]
application quadratique, dont la valeur en $z$ serait, essentiellement, $z\bar z$, mais ``\unsure{remonté|remontée}'' à $\mathcal A$.
L'application bil. associée serait
\[
(u,v)\longmapsto u\bar v+v\bar u=\operatorname{Tr}_N(u\bar v)
\]
déjà connue.

Donc il faut se donner les \unsure{autres} valeurs de $N$, et de $\varphi_N$, sur les éléments de la base $1,U_1,U_2,U_1\otimes U_2$ de $\mathcal A_1\otimes\mathcal A_2$. On devrait avoir
\[
\begin{array}{c}
N_{\mathcal A}(x_1\otimes x_2)=N_1(x_1)N(x_2)\in\Gamma\mathcal O\subset\Gamma\mathcal A\\
\Vert\\
(x_1\otimes x_2)(\bar x_1\otimes\bar x_2)\\
\Vert\\
(x_1\bar x_1)\otimes(x_2\bar x_2)
\end{array}
\]
On \unsure{peut écrire}
\[
N(1)=1,\quad N(U_1)=N_1(U_1)=c_1,\quad N(U_2)=N_2(U_2)=c_2
\]
\[
N(U_1\otimes U_2)=c_1c_2
\]
\mspage{102}
